\section*{Chapter Summary}
Latency-sensitive services running on the Java Virtual Machine (JVM) keep response times low with two mechanisms: low-pause garbage collection (GC), which collects the heap incrementally to bound pauses, and large result caches, which avoid recomputation.
These caches consume scarce DRAM and increase GC overhead, since cached objects stay reachable and the garbage collector must repeatedly scan them.
Placing them off-heap on fast storage relieves both, but introduces serialization and deserialization (S/D) on every access, raising latency.
Prior work avoids S/D cost and reduces GC scans with a storage-backed heap, but targets throughput-oriented collectors, whose long pauses are unsuitable for latency-sensitive services.

We present \chipperName{}, the first storage-backed heap for low-pause garbage collectors.
\chipperName{} extends the JVM with a second, high-capacity heap memory-mapped over an NVMe SSD, and migrates cache objects to it.
Migrated objects remain accessible as ordinary objects and are excluded from GC scans.
To preserve low latency, \chipperName{} (1) incorporates the cost of migration to the storage heap into the collector's cost model for incremental GCs, (2) integrates migration into the collector's parallel object-copying path, and (3) introduces a separate path for migrating objects that cannot be moved during incremental GCs.
We implement \chipperName{} in G1, the default garbage collector of OpenJDK, and evaluate it across four classes of latency-sensitive systems: search engines with OpenSearch, search libraries with Lucene as a foundation for many modern search engines, key-value stores with HBase, and stream-processing frameworks with Spark.
Our evaluation shows that \chipperName{} achieves up to $161\times$ lower P99 latency than DRAM-limited on-heap caches, $34\times$ lower than storage-backed off-heap caches, and $9\times$ lower than TeraHeap, a state-of-the-art storage-backed heap.
